\begin{lemma}[Deterministic one-row preprocessing map]\label{lem:step-012-one-row-map}
Under Assumption~\ref{assump:candidate-regime}, the contradiction hypothesis~\eqref{eq:local-negation},
and Proposition~\ref{prop:step-010-row-construction} and
Lemma~\ref{lem:step-010-one-use}, fix any realization \(\omega\) of
the input-independent preprocessing seed and any arbitrary
\(z,z'\in\mathsf Z_M\) with \(z\simeq z'\).  Then exactly one of the
following holds.

\begin{enumerate}
\item If \(U(\omega)>M\), the preprocessing control flow evaluates no
   coordinate of either input, so its output branch is identical on
   \(z,z'\).
\item If \(U(\omega)\le M\), the two exact size-\(n\) datasets
   \(\mathcal P_\omega(z)\) and \(\mathcal P_\omega(z')\) are equal or
   replacement-adjacent.  More precisely, if the inputs differ only at
   row \(q\), then they are equal when \(q>U\), while for \(q\le U\)
   they can differ only at the single global row \(r_q\), the \(q\)-th
   hidden-tag occurrence.
\end{enumerate}

This assertion includes \(U=0\), a changed unused record, a changed used
record, simultaneous replacement of the record's feature and label, and
arbitrary possibly nonrealizable label patterns.  All nonhidden global
rows and all hidden rows other than \(r_q\) coincide.  In particular, the
map never produces group-adjacent datasets.
\end{lemma}

\begin{proof}
Proposition~\ref{prop:step-010-row-construction} gives

\[
  U=\sum_{r=1}^n\mathbf1\{I_r=J\},
\tag{10}
\]

as a function only of \(J,I_{1:n}\), all of which belong to \(\omega\).
Hence the same value of \(U\), and therefore the same overflow branch, is
used on \(z\) and \(z'\).  When \(U>M\), Proposition~\ref{prop:step-010-row-construction} stops before either input
is read.  This proves the first branch pathwise, not only in distribution.

Suppose \(U\le M\).  Lemma~\ref{lem:step-010-one-use} gives distinct
hidden positions

\[
  1\le r_1<\cdots<r_U\le n
\tag{11}
\]

depending only on \(\omega\), and its exact row rule gives, for
\(z_\ell=(x_\ell,y_\ell)\),

\[
  \bigl(\mathcal P_\omega(z)\bigr)_{r_\ell}
  =((J,x_\ell),y_\ell),
  \qquad \ell=1,\ldots,U.
\tag{12}
\]

Every other position is a nonhidden row already fixed by \(\omega\).
The same formulas, with primes, hold for \(z'\), with exactly the same
positions \(r_\ell\) and exactly the same nonhidden rows.

If \(z=z'\), the synthesized datasets are equal.  Otherwise adjacency in
\(\mathsf Z_M\) supplies one index \(q\in[M]\) such that
\(z_\ell=z'_\ell\) for every \(\ell\ne q\).  If \(q>U\), Lemma~\ref{lem:step-010-one-use} says row \(q\) is unused, and (12)
shows equality of every used hidden row; all nonhidden rows also coincide.
Thus

\[
  \mathcal P_\omega(z)=\mathcal P_\omega(z').
\tag{13}
\]

If \(q\le U\), then for every hidden row \(r_\ell\) with \(\ell\ne q\),
the equality \(z_\ell=z'_\ell\) and (12) give equal rows.  At the sole
remaining hidden position,

\[
  \bigl(\mathcal P_\omega(z)\bigr)_{r_q}=((J,x_q),y_q),
  \qquad
  \bigl(\mathcal P_\omega(z')\bigr)_{r_q}=((J,x'_q),y'_q).
\tag{14}
\]

Equation (14) is one replacement even if both \(x_q\ne x'_q\) and
\(y_q\ne y'_q\).  It is still a legal replacement when either raw label
is inconsistent with every threshold because both labels remain elements
of \(\{0,1\}\), and Assumption~\ref{assump:central-dp} is stated on all
labeled datasets.  Every nonhidden row is a component of the fixed seed,
so no nonhidden row changes.  Thus the synthesized datasets differ in at
most the one row \(r_q\), proving replacement adjacency and excluding any
group-adjacency charge.

Finally, if \(U=0\), there are no positions in (11), every synthesized
row is nonhidden and seed-fixed, and (13) holds for every pair of inputs.
This also proves the requested null case. \end{proof}

\begin{proposition}[Fixed-seed privacy kernel with learner coins retained]\label{prop:step-012-fixed-seed-kernel}
Under Assumptions~\ref{assump:candidate-regime} and
\ref{assump:central-dp}, the contradiction hypothesis~\eqref{eq:local-negation}, Proposition~\ref{prop:step-010-simulator}, and
Lemma~\ref{lem:step-012-one-row-map}, fix a preprocessing seed \(\omega\)
but do not fix the internal coins of \(A\).  Define the all-zero global
hypothesis

\[
  h_0(i,x):=0
  \qquad ((i,x)\in X_{k,N}),
\tag{15}
\]

and the global-output kernel \(K_\omega:\mathsf Z_M\to\mathcal H_{k,N}\)
by

\[
  K_\omega(z,F)=
  \begin{cases}
    \mathbf1\{h_0\in F\}, & U(\omega)>M,\\[1mm]
    \Pr_A\!\left[A(\mathcal P_\omega(z))\in F\right],
      & U(\omega)\le M,
  \end{cases}
\tag{16}
\]

for every \(F\subseteq\mathcal H_{k,N}\).  Then, for all arbitrary
\(z\simeq z'\) in \(\mathsf Z_M\),

\[
  K_\omega(z,F)
  \le e^\varepsilon K_\omega(z',F)+\delta
  \qquad(F\subseteq\mathcal H_{k,N}).
\tag{17}
\]

On overflow the two output laws coincide exactly.  On nonoverflow, (17)
uses Assumption~\ref{assump:central-dp} at most once.  The statement holds
for arbitrary randomized \(A\); its internal coins are precisely the
randomness under \(\Pr_A\) in (16), not part of the fixed seed.
\end{proposition}

\begin{proof}
The function \(h_0\) belongs to the full output space
\(\mathcal H_{k,N}\), and \(D_Jh_0=g_0\) for every \(J\).  Thus the first
line of (16) is an analytical global-output lift of the simulator's actual
fixed overflow output; it does not call \(A\).

If \(U(\omega)>M\), the seed fixes the same branch for both inputs, and

\[
  K_\omega(z,F)=\mathbf1\{h_0\in F\}=K_\omega(z',F).
\tag{18}
\]

Since \(e^\varepsilon\ge1\) and \(\delta\ge0\), (18) implies (17).  This
also shows directly that overflow contributes no privacy residual.

Suppose \(U(\omega)\le M\).  By
Lemma~\ref{lem:step-012-one-row-map},
\(\mathcal P_\omega(z)=\mathcal P_\omega(z')\) or
\(\mathcal P_\omega(z)\simeq\mathcal P_\omega(z')\).  In the equality
case the two laws in (16) coincide, including when \(U=0\) or the changed
input row is unused.  In the adjacent case, Assumption~\ref{assump:central-dp}
applied once to the exact size-\(n\) pair and the event \(F\) gives

\[
\begin{aligned}
  K_\omega(z,F)
  &=\Pr_A[A(\mathcal P_\omega(z))\in F]\\
  &\le e^\varepsilon
    \Pr_A[A(\mathcal P_\omega(z'))\in F]+\delta\\
  &=e^\varepsilon K_\omega(z',F)+\delta.
\end{aligned}
\tag{19}
\]

No coin of \(A\) was conditioned on or fixed in (19).  Arbitrary feature
and label replacements are covered because the premise of (19) is the
replacement adjacency proved in Lemma~\ref{lem:step-012-one-row-map}, not
realizability.  There is one invocation of (7), so neither group privacy
nor sequential composition appears. \end{proof}

\begin{lemma}[Common input-independent mixture and restriction preserve one charge]\label{lem:step-012-postprocess-mixture}
Under Assumptions~\ref{assump:candidate-regime} and
\ref{assump:central-dp}, the contradiction hypothesis~\eqref{eq:local-negation}, Lemma~\ref{lem:step-010-public-preprocessing} and
Proposition~\ref{prop:step-010-simulator}, and
Proposition~\ref{prop:step-012-fixed-seed-kernel}, let \(\nu\) be the law
of the input-independent preprocessing seed on its finite seed space
\(\mathsf\Omega_{\mathrm{pre}}\).  For
\(G\subseteq\mathsf\Omega_{\mathrm{pre}}\times\mathcal H_{k,N}\), define
its seed section

\[
  G_\omega:=\{h\in\mathcal H_{k,N}:(\omega,h)\in G\},
\tag{20}
\]

and form the joint seed/global-output kernel

\[
  \widehat K(z,G)
  :=\sum_{\omega\in\mathsf\Omega_{\mathrm{pre}}}
       \nu(\omega)K_\omega(z,G_\omega).
\tag{21}
\]

Then \(\widehat K\) is \((\varepsilon,\delta)\)-DP.  Applying only after
this common mixture the deterministic postprocessing map

\[
  \Phi(\omega,h):=D_{J(\omega)}h
\tag{22}
\]

produces exactly the simulator law.  Consequently (4) holds with
exactly one additive \(\delta\).
\end{lemma}

\begin{proof}
The seed space is finite: Lemma~\ref{lem:step-010-public-preprocessing}
samples finitely many tags,
finitely many elements of the finite-support prior, and features from the
finite domain \([N]\).  For each \(G\), every section \(G_\omega\) in
(20) is therefore a valid event in the global output space.  Equation
(21) is a total kernel: all summands are nonnegative, its value on
\(\mathsf\Omega_{\mathrm{pre}}\times\mathcal H_{k,N}\) is
\(\sum_\omega\nu(\omega)K_\omega(z,\mathcal H_{k,N})=1\), and finite
additivity follows by taking disjoint seed sections and using additivity of
each \(K_\omega\).  The seedwise inequality (17) is uniform in \(\omega\)
and in the event, so for arbitrary \(z\simeq z'\),

\[
\begin{aligned}
  \widehat K(z,G)
  &=\sum_\omega\nu(\omega)K_\omega(z,G_\omega)\\
  &\le\sum_\omega\nu(\omega)
       \bigl(e^\varepsilon K_\omega(z',G_\omega)+\delta\bigr)\\
  &=e^\varepsilon\sum_\omega\nu(\omega)K_\omega(z',G_\omega)
    +\delta\sum_\omega\nu(\omega)\\
  &=e^\varepsilon\widehat K(z',G)+\delta.
\end{aligned}
\tag{23}
\]

Here Lemma \ref{lem:step-010-public-preprocessing} supplies the
same input-independent weights \(\nu(\omega)\) for \(z\) and \(z'\).  Also,
\[
\sum_\omega\nu(\omega)=1.
\]
Thus (23) first proves exact DP for the
common joint mixture.  The overflow subset
\(\{\omega:U(\omega)>M\}\), the \(U=0\) subset, and every nonoverflow
subset are determined solely by this common seed.  No input-dependent
conditioning, normalization, or change of measure occurs.

Now fix \(E\subseteq\mathsf G_N\) and use the one joint-output event

\[
  G_E:=\Phi^{-1}(E)
  =\{(\omega,h):D_{J(\omega)}h\in E\}.
\tag{24}
\]

On an overflow seed, (15) gives \(\Phi(\omega,h_0)=g_0\), so that seed's
contribution to \(\widehat K(z,G_E)\) is
\(\nu(\omega)\mathbf1\{g_0\in E\}\), independently of \(z\).  On a
nonoverflow seed, its contribution is

\[
  \nu(\omega)\Pr_A\!\left[
    D_{J(\omega)}A(\mathcal P_\omega(z))\in E
  \right],
\tag{25}
\]

with all internal randomness of \(A\) retained.  Summing these two exact
branches is precisely Proposition~\ref{prop:step-010-simulator},
so

\[
  B_{\mu_{N,M},A}(z,E)=\widehat K(z,G_E).
\tag{26}
\]

Finally apply (23) to \(G=G_E\) and use (26) for both inputs.  This is the
event-preimage proof of postprocessing after the common mixture and gives
(4).  The additive term in (23) is \(\delta\), not \(\delta\) times the
number of seeds, \(n\), \(k\), \(U\), or an overflow factor. \end{proof}

\begin{proposition}[Exact source-cap membership]\label{prop:step-012-source-membership}
Under Assumptions~\ref{assump:candidate-regime} and
\ref{assump:central-dp}, the contradiction hypothesis~\eqref{eq:local-negation}, Propositions~\ref{prop:step-005-certificate} and
\ref{prop:step-010-simulator}, and
Lemma~\ref{lem:step-012-postprocess-mixture}, define the exact source-cap
learner set at the present \(N,M\) by

\[
\begin{aligned}
  \mathcal K^{\mathrm{src}}_{N,M}:=
  \bigl\{K:\;&\mathsf Z_M\longrightarrow\mathsf G_N
      \text{ a total randomized kernel}:\\[-1mm]
    &K(z,E)\le e^{0.1}K(z',E)+\Delta_M\\[-1mm]
    &\text{for all }z\simeq z'\text{ in }\mathsf Z_M
      \text{ and }E\subseteq\mathsf G_N\bigr\}.
\end{aligned}
\tag{27}
\]

Then

\[
  B_{\mu_{N,M},A}\in\mathcal K^{\mathrm{src}}_{N,M}.
\tag{28}
\]

More fully, the same \(N,M\) satisfy the source scalar conditions in
(8), the simulator has exact input size \(M\), full arbitrary one-block
output space \(\mathsf G_N\), and is total on all labeled inputs, and its
event probabilities satisfy (6).  Thus (28) is exact source-interface
membership, not an inference from the phrase
"\((\varepsilon,\delta)\)-DP" alone.
\end{proposition}

\begin{proof}
Proposition~\ref{prop:step-010-simulator} proves that
\(B_{\mu_{N,M},A}\) is a total randomized kernel with domain exactly
\(\mathsf Z_M\) and codomain exactly the full improper space
\(\mathsf G_N\).  Its totality includes overflow, \(U=0\), arbitrary
corrupt input labels, and unused trailing records.

Proposition~\ref{prop:step-005-certificate}, applied to the same
\(N,M,\varepsilon,\delta\), gives (8)--(9).  In particular, \(M\ge8\)
implies \(\log M>0\), so \(\Delta_M\) in (5) is a well-defined positive
source cap.  Lemma~\ref{lem:step-012-postprocess-mixture} gives, for every
defining pair and event in (27),

\[
  B_{\mu_{N,M},A}(z,E)
  \le e^\varepsilon B_{\mu_{N,M},A}(z',E)+\delta.
\tag{29}
\]

Because probabilities are nonnegative, \(\varepsilon\le0.1\), and
\(\delta<\Delta_M\),

\[
\begin{aligned}
  e^\varepsilon B_{\mu_{N,M},A}(z',E)+\delta
  &\le e^{0.1}B_{\mu_{N,M},A}(z',E)+\delta\\
  &< e^{0.1}B_{\mu_{N,M},A}(z',E)+\Delta_M.
\end{aligned}
\tag{30}
\]

The weak inequality required in (27) follows.  Equations (8), (27), and
(30), together with the exact total kernel typing, check every component
of the source interface separately.  No accuracy, realizability, or
properness property has been added to that learner set, and no private
eligibility is inferred without the event-level derivation (29)--(30).
This proves (28). \end{proof}
